Dues càrregues de $3,0\ \mu\mathrm{C}$ estan localitzades a $x = 0\ \mathrm{m}$, $y = 2,0\ \mathrm{m}$ i a $x = 0\ \mathrm{m}$, $y = -2,0\ \mathrm{m}$. Dues càrregues més, de valor $Q$, estan localitzades a $x = 4,0\ \mathrm{m}$, $y = 2,0\ \mathrm{m}$ i a $x = 4,0\ \mathrm{m}$, $y = -2,0\ \mathrm{m}$.

\begin{apartat}{1}
Si a l'origen de coordenades el camp elèctric és $4,0\times 10^{3}\ \mathrm{N\,C^{-1}}$ en la direcció de l'eix $x$ en sentit positiu, calculeu el valor de les càrregues.
\end{apartat}

\begin{apartat}{1}
Si el valor de les càrregues fos $Q = 2,0\ \mu\mathrm{C}$, calculeu la força $\vec{F}$ que experimentaria un protó situat a l'origen de coordenades.
\end{apartat}

\dades{%
  $k = \dfrac{1}{4\pi\varepsilon_0} = 8,99\times 10^{9}\ \mathrm{N\,m^2\,C^{-2}}$.\\
  Càrrega elemental $= 1,60\times 10^{-19}\ \mathrm{C}$.}
